\section{Analysis Metric 1: File Types}
\label{sec:analysisFileTypes}

\subsection{Metric Definition}
\label{sec:metricFileTypesDefinition}

The first metric we analyzed concerns the file types of \ac{NPM} packages.

To compute this metric, the developed tool (\cref{sec:ToolForAnalyzingNpmPackages}) uses 
\texttt{Magika}, a deep learning-based content-type detection tool developed by Google~\cite{magika}.

The developed tool first determines the type of each file present in the package version 
and then constructs a list of all unique file types found.

This metric was taken into consideration 
because the attack A described in \cref{sec:attackA} introduces a malicious Windows \ac{DLL} 
into the compromised package versions.


\subsection{Goodware Dataset Analysis}
\label{sec:analysisFileTypesGoodware}

For this metric, the Goodware dataset (explained in \cref{sec:goodwareDataset}) contains 
20 lists of file types, one for each version of every package.

We used heatmaps, \cref{fig:fig1ComFileTypes} and \cref{fig:fig1RareFileTypes},
to visualize the distribution and evolutionary patterns 
of the file types of the packages in the Goodware dataset across versions.

In the heatmaps, the x-axis shows the package versions, from the least recent (-19) to the latest available (0),
while the y-axis represents the file extensions.

The numbers inside the cells are the relative frequency
of the extensions respect to the total number of packages in the dataset.

A value like 1.000 indicates that all packages in the dataset (100\%) have at least one file of that type in that version.
Conversely, a value such as 0.010 indicates that only 1\% of the packages have that extension among their files.

Therefore, this number also indicates how common it is to find a given extension in the packages.

The heatmaps use dark color shades for high values and light shades for low values.

The \cref{fig:fig1ComFileTypes} and \cref{fig:fig1RareFileTypes} show all seventy-one file types found 
in the dataset, ordered from the most common to the rarest.

We used two heatmaps to better visualize the analysis:
\begin{enumerate}
    \item \textbf{Common File Types:} \cref{fig:fig1ComFileTypes} is the heatmap showing file types 
    that have a relative frequency higher than 1\% in at least one version.
    \item \textbf{Rare File Types:} \cref{fig:fig1RareFileTypes} is the heatmap showing all remaining types, 
    which have a relative frequency that does not exceed 1\% across the versions.
\end{enumerate}